\section{Evaluation Design}\label{sec:design}

The evaluation follows the chain, and each comparison permits one kind of statement. A construction check asks whether the adapters encode what they were trained on (anchor to instrument). A second-bank check asks whether the encoding of the anchor's order survives a change of scenario text. A factorial asks whether the declared dimensions behave as separate inputs. A held-out screen asks whether the instrument's readings of new items agree with an existing human survey (instrument to respondents). Descriptive follow-up analyses, added after the registered results had been read, place the instrument beside the respondents and beside its own training labels. Every number carries the label of the surface it was computed on (Table~\ref{tab:surfaces}), and numbers from different surfaces are never pooled.

\begin{table}[htbp]
\centering
\caption{Evaluation surfaces. Every number in the paper carries one of these labels.}\label{tab:surfaces}
\small
\renewcommand{\arraystretch}{1.15}
\begin{tabularx}{\textwidth}{@{}l>{\raggedright\arraybackslash}X>{\raggedright\arraybackslash}X>{\raggedright\arraybackslash}p{0.27\textwidth}@{}}
\toprule
Label & Units and items & Comparison & What it can show \\
\midrule
S1 & First-cohort country adapters (original bank), strict battery & Adapter--GPS rank correlation over countries; model side only & Whether a surface encodes its declared sign \\
S2 & Second-cohort profile adapters (second bank), each country mapped to its profile, strict battery & Same statistic and bootstrap as S1 & Whether that encoding survives a new scenario bank \\
H & Second-cohort adapters on held-out WVS Wave~7 items & Signed agreement of selected items with human country means, against GPS & Retrospective screen of new items against an existing survey \\
F & Second-cohort adapters and prompted profiles, strict battery & Main effects of the six sign bits on seven surfaces; model side only & Whether the declared dimensions separate \\
R & One profile retrained without a fixed seed & Score differences against the original adapter; model side only & Retraining noise for one unit \\
N & Two near-zero refits and size-matched controls & Composite shifts in profile standard deviations; model side only & Sensitivity of the sign rule's margin \\
P & Country-named persona prompt on the base model, strict battery & Persona--GPS rank correlation; paired difference with S1; model side only & Whether naming the country carries the GPS ordering \\
D & First-cohort adapters, second-cohort factorial and the persona routes, strict battery; descriptive follow-up analyses, not preregistered & Agreement with the GPS and with the WVS Wave~7 composite; distance to survey answer distributions; decomposition of main effects; held-out label recovery & A description of how the instrument relates to respondents and to its own labels; no test \\
\bottomrule
\end{tabularx}
\end{table}

\subsection{Second-bank check}\label{sec:design-bank}

For each surface of the strict battery, we compute the Spearman correlation between the adapters' surface composites and the GPS coordinate of the surface's own dimension over the \dc{42}{countries} countries, first for the first cohort (S1) and then for the second cohort with each country assigned its profile's adapter (S2). Intervals come from a profile-block bootstrap that resamples the \dc{26}{profiles_realized} realized sign profiles, with \dc{2,000}{boot_draws} draws, seed \dc{20261011}{seed_c3}, and percentile limits. The imprint set is the \dc{4}{imprint_dims} dimensions whose first-cohort interval excluded zero: trust, risk taking, positive reciprocity and altruism. Patience and negative reciprocity were flat in the first cohort and are reported without interpretation. The rule was fixed before the second cohort was trained: an imprint dimension whose second-bank interval includes zero is named, and no sentence about that dimension carrying over to the second bank is written. We also report the Pearson agreement between the two banks over the \dc{26}{profiles_realized} shared profiles, with a \dc{2,000}{boot_draws}-draw bootstrap over profiles. It measures how much the scenario text changes the instrument and involves no human data.

\subsection{The country-named prompt}\label{sec:design-persona}

The persona route (Section~\ref{sec:instrument}) was scored on the strict battery for each of the \dc{42}{countries} countries in one run on the frozen stack, after a rescoring of the base model that reproduced the banked base scores exactly. We compute the persona--GPS correlation for each surface (surface P) and one planned contrast, fixed in a design note committed before any persona was scored: the first cohort's trust correlation minus the persona route's, with both recomputed on every draw of the profile-block bootstrap of Section~\ref{sec:design-bank}. The other differences, and those for the second cohort and the sign-profile prompt, are descriptive.

\subsection{The held-out screen}\label{sec:design-screen}

\paragraph{Item pool.} The pool was built from WVS Wave~7 metadata alone: substantive ordinal questions with two to eleven contiguous response codes that lie outside the battery used to develop the adapters. A text review, with a recorded reason for every change, excluded items with non-monotone codes or factual content and repaired wording that did not extract from the questionnaire. The result is \dc{223}{pool_items} items that no adapter was built for (online supplement).

\paragraph{Probe.} All \dc{64}{cells} second-cohort adapters scored every pool item under the frozen specification. For each item $i$, the probe regresses the \dc{64}{cells} adapter scores on an intercept, the six sign bits coded $\pm\dc{0.5}{x_code}$, and their \dc{15}{interactions} two-way interactions, which leaves \dc{42}{error_df} error degrees of freedom. With $t_{d}$ the main-effect $t$ statistic of dimension $d$, the item's assigned dimension is $d_i=\arg\max_d|t_d|$, or none when the largest $|t_d|$ is below \dc{2.0}{t_threshold}, and its sign is $s_i=\operatorname{sign}(b_{d_i})$. The primary rule keeps items whose assigned dimension is in the imprint set and whose $|t|$ is at least \dc{2}{t_threshold}, up to the top \dc{25}{top_share} percent of all pool items by $|t|$, which caps the selection at \dc{56}{max_selected} items. The secondary rule keeps every imprint item above the threshold; it was declared a reliability check and does not enter the inference. The probe reads model scores only. The selection rule does not control the rate at which noise items cross the threshold, and the permutation null below absorbs it.

\paragraph{Statistic and criterion.} For each selected item, the signed item agreement is
\begin{equation}
a_i = s_i\cdot\rho\big(\bar h_{ci},\,z_{c,d_i}\big),
\qquad
A=\frac{1}{|S|}\sum_{i\in S}a_i,
\label{eq:screen}
\end{equation}
where $\rho$ is the Spearman correlation over the \dc{42}{countries} countries, $\bar h_{ci}$ is the survey-weighted WVS Wave~7 mean of item $i$ in country $c$ with negative codes set to missing, and $S$ is the selected set. An item--country cell that was never fielded enters as missing and is not imputed. $A$ is positive when the human country differences on the selected items run in the direction the instrument assigned, against the GPS coordinate of the assigned dimension.

\paragraph{Inference.} The registered test is a permutation null over items: draw as many items as were selected uniformly from the pool, give them the selected items' dimension and sign pairs in random order, and recompute $A$; \dc{10,000}{perm_draws} draws, seed \dc{20261015}{seed_screen}, one-sided $p$ equal to the share of null draws at or above the observed $A$. Because the null redraws items from the same pool, it absorbs the rate at which the probe flags noise items. A \dc{95}{ci_level} percent interval for $A$ comes from \dc{2,000}{boot_draws} bootstrap resamples of the \dc{42}{countries} countries. The registered readings were fixed in advance: if $A$ is not above the null at one-sided $\alpha=\dc{0.05}{alpha}$, the adapters do not screen held-out items better than random assignment; if $A$ is above the null but below one third of the human ceiling, a screening signal exists but is weak.

\paragraph{Human ceiling.} For each selected item fielded in both WVS waves 6 and 7 we correlate the two waves' country means over the \dc{28}{ceiling_countries} countries with both waves \citep{wvsts2022} and average over items, which gives the agreement two human surveys reach on these items. The ratio of $A$ to this ceiling is a benchmark only, because the two statistics are built differently.

\paragraph{Per-dimension readings.} The registered analysis also reports $A_d$, the mean of $a_i$ over the selected items assigned to dimension $d$, with an interval from the same country resamples. An interval that includes zero gives no evidence of agreement for that dimension. These intervals are unadjusted for the four imprint dimensions, so we add exploratory Bonferroni intervals read off the same draws at the \dc{98.75}{bonf_level} percent level.

\paragraph{Comparators.} Two comparators read only the questionnaire and saw no model score or human response: a face-valid rule that assigns an item to the one dimension whose construct words its wording contains, and a typed decision model (jev-1.13.0) run \dc{5}{jev_runs} times, which assigns an item when at least \dc{3}{jev_majority} runs agree on dimension and sign (online supplement). The registered primary contrast is coverage-matched, $D_{\mathrm{common}}=A_{\mathrm{adapter}}-A_{\mathrm{comparator}}$ on the items both assign, with an interval from the paired country bootstrap. An interval above zero would mean the adapters find structure the wording does not reveal; one that covers zero was registered as ``the adapters match a careful reader''; one below zero would favor the reader. The contrast on each route's own items is reported beside it. The preregistration named in advance that both comparators concentrate on trust.

\paragraph{Readout variants and the prompt route.} The probe was re-estimated on three re-scorings of the pool (options reversed, and two paraphrases of the prompt template); the selection was declared dependent on wording if more than half of the selected items changed dimension or sign. The same probe and test were applied to the \dc{64}{cells} prompted sign profiles.

\subsection{Registration, blinding and amendments}\label{sec:design-prereg}

The screen was registered in a document committed to the project repository and timestamped by a pushed git tag before any held-out item was scored by an adapter; at that point no selection existed and no held-out human value had been read. A freeze guard enforced the order of work: score the held-out items, run the probe, freeze the hashes of the selection and the item pool, and only then join the human data.

Two amendments were made, both before any human held-out value was read (Figure~\ref{fig:timeline}). Amendment~1 repaired \dc{12}{repaired_items} items whose standalone wording could not be answered as written, under the pool rule and override mechanism already registered, and every route was re-scored as one run. We disclose that the first probe pass had been computed and read before the repairs were chosen, and that two repaired items (Q153 and Q157) were in its top quartile. The comparators had assigned the pool on the pre-repair wording, so Amendment~2 removed the repaired items from both comparator routes before any comparator statistic was computed; it affected one item (Q82), which was not in the adapter selection. That decision preceded any human value, and the amendment text, written after the primary statistic was read, records it.

The criterion is retrospective. The WVS Wave~7 file is public, existed before the probe, the pool and the adapters, and is the only human source in this study. Registration guards against tuning the probe to the criterion; it cannot make the criterion prospective. Results on surface H are therefore labelled retrospective throughout. The online supplement discloses the provenance gaps: two code revisions in one run, an unnamed commit for the first cohort's scores, and an unexplained gap to the development-stage archive, whose scores no statistic here uses.

\begin{figure}[htbp]
\centering
\begin{tikzpicture}[>=Latex,node distance=4mm,every node/.style={font=\footnotesize,align=center},
  wstep/.style={draw,rounded corners,text width=0.78\textwidth,minimum height=7mm},
  human/.style={wstep,fill=black!12,line width=0.9pt}]
\node[wstep] (pool) {Item pool built from WVS Wave~7 metadata and questionnaire text only};
\node[wstep,below=of pool] (tag) {Preregistration committed and timestamped by a pushed git tag};
\node[wstep,below=of tag] (s1) {First scoring pass of the pool, adapters only; probe read};
\node[wstep,below=of s1] (a1) {Amendment 1: \dc{12}{repaired_items} wording repairs; every route re-scored as one run};
\node[wstep,below=of a1] (sel) {Probe selection on the amended scores (model side only)};
\node[wstep,below=of sel] (a2) {Amendment 2 decided: repaired items leave both comparator routes; freeze file written};
\node[human,below=of a2] (join) {First read of held-out human data (retrospective WVS Wave~7 criterion)};
\node[wstep,below=of join] (exp) {Registered analyses, then labelled exploratory analyses};
\foreach \a/\b in {pool/tag,tag/s1,s1/a1,a1/sel,sel/a2,a2/join,join/exp} \draw[->] (\a)--(\b);
\end{tikzpicture}
\caption{Order of work for the held-out screen. Every step above the shaded box used model scores, metadata or questionnaire text only. The shaded box is the single read of held-out human data.}\label{fig:timeline}
\end{figure}

\subsection{The sign-profile factorial}\label{sec:design-factorial}

The second cohort covers all \dc{64}{cells} sign profiles, so the strict-battery composites form a full factorial in six binary factors. For each of the seven surfaces we fit the probe's regression (intercept, six sign bits coded $\pm\dc{0.5}{x_code}$, \dc{15}{interactions} two-way interactions, \dc{42}{error_df} error degrees of freedom). A main effect is the difference between the mean composites of profiles with the bit set and without it. Because the bits were assigned by design, this is a causal statement about the instrument, which acts on synthetic labels; it carries no implication about how preferences operate in people. Each bit has an own surface (patience owns two). The project's analysis plan set one acceptance check, separate from the held-out registration: at least \dc{4}{v5_required} of the six bits must move their own surface more than any other (V5), against a binomial reference with a one-in-seven chance of a hit. Exploratory scale-free views divide effects by each surface's standard deviation over the \dc{64}{cells} profiles, count the surfaces whose largest mover is their own bit, and report the first principal component's variance share and the mean absolute correlation between surfaces. The same analyses on the \dc{64}{cells} prompted profiles contrast two instruments that carry the same declared information.

\subsection{Retraining noise and the near-zero margin}\label{sec:design-noise}

One profile adapter was retrained under the same protocol without a fixed seed, which estimates a noise floor for the pipeline; we compare its scores with the original on both batteries. A predeclared sensitivity check also removes the training pairs of every dimension whose GPS coordinate lies within \dc{0.10}{near_zero_cut} of zero (the sign rule treats a coordinate just above zero like a large positive one) and refits two countries (Great Britain and Colombia), each beside a size-matched control. The contrast is refit minus control on the dimensions still in training, in profile standard deviations, read against the retrain's shift. The rule was to report the margin beside the main results and never to promote it.

\subsection{Descriptive follow-up analyses}\label{sec:design-desc}

Four further analyses were added after the registered results had been read. They are not preregistered, enter no registered family, and test no rule fixed in advance; we label their numbers D and read them as descriptions. Their human input is the WVS Wave~7 file used for the screen, so every comparison with respondents is retrospective. All use the strict battery and the \dc{42}{countries} countries.

\paragraph{Three columns.} For each of the seven surfaces we compute three Spearman correlations over countries: the adapter composite with the GPS coordinate of the surface's own dimension, the adapter composite with the Wave~7 composite of the same items, and that human composite with the GPS. The adapter columns use the first-cohort adapters (S1; the second cohort is in the supplement) with the profile-block bootstrap of Section~\ref{sec:design-bank}; the human--GPS interval resamples countries (seed \dc{20260905}{seed_country}). The human columns of risk taking cover one country fewer, because one country has no valid answer to the risk item.

\paragraph{Distance to the survey distributions.} For country $c$ and item $q$, the total variation distance between the model's option probabilities and the survey-weighted human answer shares is
\begin{equation*}
\mathrm{TVD}_{cq}=\tfrac{1}{2}\sum_j \bigl|\widehat{P}_c(j\mid x_q)-P^{H}_c(j\mid q)\bigr|,
\end{equation*}
where $j$ runs over the answer codes and $P^{H}_c$ is the share of valid human answers equal to $j$. The headline statistic is the mean over items within a country and then over countries, so that each country counts once. We score four cells on the same frozen stack: the base model, the first-cohort adapters, the country-named persona prompt on the base model (surface P), and the adapters with the persona prompt together (gates in the supplement). Contrasts are paired differences of country means. The interaction contrast is the adapters' effect with the persona prompt minus their effect without it, and zero means that the two add. Each contrast has a profile-block and a country interval (\dc{2,000}{boot_draws} draws each). The human answer shares enter the distances and are not reported.

\paragraph{Cross-talk decomposition.} Starting from the main effects of the second-cohort factorial (Section~\ref{sec:design-factorial}), we write the main-effects prediction of a surface at a country's sign profile as the sum over the six bits of each bit's main effect times its code, and split it into an own part (the surface's own bit) and a cross part (the other five). We correlate each part with the GPS coordinate of the surface's own dimension over the countries and report Var(cross)/Var(full). The decomposition is behavioral and uses model quantities and published GPS values only.

\paragraph{Held-out label recovery.} Training pairs were split once into a training part and a held-out part that the adapters never saw. For each first-cohort adapter we scored each held-out pair as the sequence log probability of the chosen and of the rejected completion, under the adapter and under the base model, in a further run on the same stack (gates in the supplement). Recovery, accuracy and base accuracy are defined in Section~\ref{sec:res-recovery}. This is a construction check on the synthetic labels and involves no human data.